\section{An exact endpoint criterion}
\label{sec:endpoint-criterion}
The nonsummable tail has a precise role in the construction: the endpoint
comparisons force a dual label that the actual dual space cannot contain.
We now separate this requirement from the particular triangular coordinates.
The following theorem gives a criterion within a specified class of graphs;
it does not assume that arbitrary maximally monotone operators admit this
representation.

Let $X$ be a real Banach space and let $H_0$ be a real Hilbert space.
In this section put $H=\R\oplus H_0$. Let
$E:X^*\to X$ and $M:X^*\to H$ be bounded linear maps, and let
$g\in X^*\setminus\{0\}$ satisfy $Mg=0$ and
\begin{equation}
\pair{Ea}{b}+\pair{Eb}{a}=2(Ma,Mb)_H
\qquad(a,b\in X^*).
\label{eq:abs-bilinear}
\end{equation}
Define $h=Eg$, $r(a)=\pair{h}{a}$, $La=-Ea+r(a)h$, and
$K=\ker M\cap\ker r$. Assume the exact primal annihilator identity
\begin{equation}
\{x\in X:\pair{x}{k}=0\ \hbox{for all }k\in K\}=\R h.
\label{eq:abs-annihilator}
\end{equation}
Let $\omega:(0,\infty)\to H_0$ be $\ell$-Lipschitz, where
$0\le\ell\le1$, and suppose
\begin{equation}
\omega(s)\longrightarrow\eta\text{ in }H_0\quad(s\downarrow0),
\qquad 0<S:=\|\eta\|^2<1,\qquad \|\omega(s)\|^2\le S.
\label{eq:abs-tail}
\end{equation}
Put $J=(-\infty,-1)\cup(1,\infty)$ and
$F(t)=(\sgn t,\omega(|t|-1))$. For every $t\in J$, assume
there exists an actual $a^t\in X^*$ with $r(a^t)=t$ and
$Ma^t=F(t)$. We use every label with these values, defining
\begin{equation}
D=\{a\in X^*:r(a)\in J,\ Ma=F(r(a))\},\qquad
G=\{(La,a):a\in D\}.
\label{eq:abs-full-graph}
\end{equation}
In particular, each parameter fibre is $a^t+K$, without a boundedness
or continuity assumption on the choice $t\mapsto a^t$.

\renewcommand{\thetheorem}{E}
\begin{theorem}[Actual endpoint labels and maximality]
\label{thm:endpoint-criterion}
Under the preceding assumptions, define
\[
\widehat F(t)=
\begin{cases}
F(t),&|t|>1,\\
(t,\eta),&|t|\le1.
\end{cases}
\]
Then the entire monotone polar in $X\times X^*$ is
\begin{equation}
G^\mu=\{(Lp,p):p\in X^*,\ Mp=\widehat F(r(p))\}.
\label{eq:abs-entire-polar}
\end{equation}
It is the unique maximally monotone extension of $G$. Moreover, $G$
itself is maximally monotone if and only if
\begin{equation}
\nexists(p,\tau)\in X^*\times[-1,1]:
\qquad r(p)=\tau,\quad Mp=(\tau,\eta).
\label{eq:abs-endpoint-exclusion}
\end{equation}
When Eq.~\eqref{eq:abs-endpoint-exclusion} holds, the operator
$\mathcal A$ with graph $G$ omits zero and satisfies
$\mathcal V(G)\le S\|g\|$. With
$\mathcal P x=\pair{x}{g}g$, both $\mathcal A$ and $\mathcal P$
are maximally monotone, satisfy
$\dom\mathcal A\cap\operatorname{int}\dom\mathcal P\ne\varnothing$,
and have a nonmaximal sum: $(0,0)$ belongs to its monotone polar
and lies outside its graph.
\end{theorem}
\begin{proof}
Equation~\eqref{eq:abs-bilinear} at $a=b=g$ gives $r(g)=0$.
Using $Mg=0$ again gives
\begin{equation}
\pair{Ea}{g}=-r(a),\quad \pair{La}{g}=r(a),\quad
\pair{La-Lb}{a-b}=|r(a)-r(b)|^2-\|Ma-Mb\|^2.
\label{eq:abs-pairing}
\end{equation}
The actual label at $t=2$ also proves $h\ne0$.

Extend $\omega$ to zero by $\omega(0)=\eta$; its Lipschitz bound
persists by the norm limit. Define
$\chi(t)=\max\{-1,\min\{1,t\}\}$ and
$\zeta(t)=\max\{|t|-1,0\}$. Then
$\widehat F(t)=(\chi(t),\omega(\zeta(t)))$.
If $t,u$ lie on the same exterior branch, in the middle interval,
or one in the middle and one outside, the nonnegative changes in
$\chi$ and $\zeta$ add to $|t-u|$. On opposite exterior branches
they sum to $2+||t|-|u||\le|t-u|$. Thus in all cases
\[
|\chi(t)-\chi(u)|+|\zeta(t)-\zeta(u)|\le|t-u|,
\]
and
\[
\|\widehat F(t)-\widehat F(u)\|^2
\le|\chi(t)-\chi(u)|^2+\ell^2|\zeta(t)-\zeta(u)|^2
\le|t-u|^2.
\]
Equation~\eqref{eq:abs-pairing} proves monotonicity of both $G$
and the set on the right of Eq.~\eqref{eq:abs-entire-polar}.
It also proves that this latter set is contained in $G^\mu$.

For the reverse inclusion, take an arbitrary $(x,p)\in G^\mu$,
where $p$ belongs to the full continuous dual. For $a\in D$ and
$t=r(a)$, expansion using Eq.~\eqref{eq:abs-bilinear} gives
\begin{align*}
\pair{x-La}{p-a}
={}&\pair{x}{p}-\pair{x+Ep}{a}
       +2(Ma,Mp)-t r(p)+t^2-\|Ma\|^2.
\end{align*}
Replacing $a$ by $a+\theta k$, for all $\theta\in\R$ and
$k\in K$, forces $\pair{x+Ep}{k}=0$. Hence
$x=-Ep+v h$ by Eq.~\eqref{eq:abs-annihilator}. With
\[
\tau=\frac{v+r(p)}2,\qquad \nu=\frac{v-r(p)}2,
\]
the same expansion reduces exactly to
\begin{equation}
\pair{x-La}{p-a}=(t-\tau)^2-\nu^2-\|F(t)-Mp\|^2.
\label{eq:abs-polar-square}
\end{equation}
Every $t\in J$ is an actual test. If $|\tau|>1$, testing
$t=\tau$ forces $\nu=0$ and $Mp=F(\tau)$; thus
$v=r(p)=\tau$ and $x=Lp$.

If $|\tau|\le1$, write $Mp=(b,z)\in\R\oplus H_0$.
Limits along the two actual branches yield
\[
(b-1)^2+\|z-\eta\|^2+\nu^2\le(1-\tau)^2,
\qquad
(b+1)^2+\|z-\eta\|^2+\nu^2\le(1+\tau)^2.
\]
Multiply these by $(1+\tau)/2$ and $(1-\tau)/2$, respectively,
and add. Expansion gives
\[
(b-\tau)^2+\|z-\eta\|^2+\nu^2\le0.
\]
The endpoint weights remain valid at $\tau=\pm1$.
Consequently $Mp=(\tau,\eta)$, $r(p)=v=\tau$, and $x=Lp$.
This proves Eq.~\eqref{eq:abs-entire-polar}. The norm limits were
taken in $H$, not in $X^*$; the label $p$ was actual from the start.

We have proved that $G^\mu$ is monotone. Any point related to it
is related to $G$ and therefore belongs to $G^\mu$, so $G^\mu$
is maximally monotone. Every monotone extension of $G$ lies in
$G^\mu$; hence a maximal one must equal it. The missing part
$G^\mu\setminus G$ is exactly the set of labels in
Eq.~\eqref{eq:abs-endpoint-exclusion}. This proves the criterion.

Finally, each $a\in D$ gives
\[
|\pair{La}{g}|=|r(a)|>1,\quad
\|La\|>1/\|g\|,\quad
\pair{La}{a}=r(a)^2-1-\|\omega(|r(a)|-1)\|^2>-S.
\]
These imply the missing zero and the stated whole-graph radial bound.
The map $\mathcal P$ is bounded, linear, positive, nonzero, and full-domain.
For an arbitrary point $(z,z^*)$ in its polar, tests at $z+tv$ give
\[
-t\pair{v}{z^*-\mathcal Pz}+t^2\pair{v}{g}^2\ge0
\qquad(t\in\R,\ v\in X).
\]
Both signs of arbitrarily small $t$ force $z^*=\mathcal Pz$.
Thus $\mathcal P$ is maximally monotone, and the actual label at
$t=2$ gives the ordered intersection condition. For every sum row,
\[
\pair{La}{a+\mathcal PLa}
=2r(a)^2-1-\|\omega(|r(a)|-1)\|^2>1-S>0.
\]
The sum is monotone and omits the zero primal. Adjoining $(0,0)$
is a proper monotone extension.
\end{proof}

For the concrete example, use the maps $E,m,g$ of
Eqs.~\eqref{eq:D4}--\eqref{eq:D5}, $H_0=\ell^2(\N)$, and
$\omega(s)=\omega_s$. Equation~\eqref{eq:C1} supplies the
bilinear identity. The complete-kernel calculation in the proof of
Theorem~\ref{thm:C} supplies the primal annihilator identity,
and Eq.~\eqref{eq:D9} supplies every actual parameter fibre.
The tail bound and limit give $S=\sigma$. Since every $mp$ is
summable and $\eta=(1/(4n))_{n\ge1}$ is not, the endpoint
exclusion holds. Thus the theorem explains which properties of the
coordinate construction are needed and why the Lipschitz extension
$\widehat F$ in the Hilbert comparison space does not add a graph point.

\noindent\textit{Reverse design from the polar condition.}
The criterion gives a way to reconstruct the design requirements from
the full polar: retain all kernel fibres so that every polar query has
the scalar form in Eq.~\eqref{eq:abs-polar-square}, realize the two
exterior branches, and choose a common limiting tail with no actual
dual label. The bound in Eq.~\eqref{eq:abs-tail} then supplies the
finite radial estimate and the rank-one partner. For the present model,
the gap between $\ell^1$ and $\ell^2$ provides exactly that absent
label. This is a mathematical reconstruction of the mechanism, rather
than an additional assumption or a replacement for the preceding
coordinate proof.
